\chapter{Austerity, Scarcity, and Fiscal Morality}
\label{ch:austerity}

\epigraph{The national budget must be balanced. The public debt must be reduced. The arrogance of the authorities must be moderated and controlled.}{Attributed to Cicero (apocryphal)}

\noindent
Austerity narratives---the claim that public spending must be cut, that governments must ``live within their means,'' that debt is inherently dangerous---are among the most successful doxonomic products of the past century.
This chapter analyzes their construction and effects.

\section{The Household Metaphor}
\label{sec:ch25-household}

\begin{definition}[Fiscal Morality Narrative]
\label{def:fiscal-morality}
\index{fiscal morality}
The \emph{fiscal morality narrative} is the claim that sovereign budgets are analogous to household budgets: governments, like families, must not spend more than they earn, must avoid debt, and must ``tighten their belts'' in hard times.
\end{definition}

This analogy is economically misleading---sovereign currency issuers face fundamentally different constraints than households---but doxonomically powerful because it maps complex macroeconomics onto intuitive domestic experience.

\section{Crisis as Narrative Opportunity}
\label{sec:ch25-crisis}

\begin{model}[Shock Doctrine]
\label{mod:shock-doctrine}
\index{shock doctrine}
Following \textcite{klein2007}, economic crises create windows of doxonomic opportunity.
During a crisis, public attention is high, uncertainty is high, and trust in existing narratives is shaken.
The first coherent narrative to fill the vacuum captures the frame.
Formally, let $H(t)$ be narrative entropy at time $t$.
A crisis causes $H(t) \to H_{\text{max}}$ (maximum uncertainty), followed by rapid entropy reduction as a new dominant frame is installed:
\[
  H(t) = H_{\text{max}} \cdot e^{-\lambda(t - t_{\text{crisis}})} + H_{\text{new}}
\]
The group that supplies the new frame controls the post-crisis common sense.
\end{model}

\section{The 2008--2013 Austerity Turn: A Doxonomic Case Study}
\label{sec:ch25-case}

The aftermath of the 2008 financial crisis is one of the best-documented examples of rapid doxonomic reframing.

\textbf{Phase 1: Crisis narrative (2008--2009).}
The initial frame was ``banks caused this.''
Public anger targeted financial institutions.
Narrative entropy was high.
Government intervention (bailouts, stimulus) was initially popular.

\textbf{Phase 2: Reframing (2009--2010).}
Within 18 months, the dominant frame shifted to ``governments spent too much.''
This reframing was not spontaneous.
It was driven by:
\begin{itemize}[nosep]
  \item think tanks (e.g., the Reinhart-Rogoff 90\% debt threshold, later shown to contain spreadsheet errors, was heavily promoted before peer review),
  \item financial press framing sovereign debt as the ``real crisis,''
  \item European institutions (ECB, European Commission) conditioning bailout assistance on austerity measures,
  \item political actors who saw an opportunity to advance pre-existing preferences for smaller government.
\end{itemize}

\textbf{Phase 3: Normalization (2010--2013).}
Austerity became common sense across much of Europe and influenced U.S.\ fiscal policy.
The household metaphor (``belt-tightening'') dominated media framing.
Counter-narratives (Keynesian stimulus, MMT) existed but lacked comparable institutional infrastructure.

\begin{remark}
The austerity turn illustrates several doxonomic mechanisms simultaneously:
the shock doctrine model (crisis opens the Overton window),
narrative infrastructure effects (think tanks had pre-packaged arguments ready for deployment),
credibility laundering (an academic paper gave the 90\% threshold a veneer of scientific authority),
and common-sense capture (within two years, austerity went from a policy choice to ``the only responsible option'').
\end{remark}

\section{Negative Cases}
\label{sec:ch25-negative}

Not all austerity narratives succeed.
Iceland rejected bank bailouts and austerity in 2008--2009 through a popular referendum; its recovery was faster than many austerity-adopting countries.
Several Latin American countries in the 2000s explicitly rejected IMF austerity programs and experienced growth.
These cases demonstrate that the austerity narrative requires specific institutional conditions to achieve capture---it is not inevitable even in a crisis.

\section{Notes and References}
\label{sec:ch25-notes}

On austerity as ideology, see \textcite{streeck2016} and \textcite{brown2015}.
The shock doctrine is developed in \textcite{klein2007}.
The political economy of debt narratives connects to \textcite{graeber2011}.
For the history of neoliberal fiscal narratives, see \textcite{mirowski2009}, \textcite{harvey2005}, and \textcite{slobodian2018}.
On the Reinhart-Rogoff controversy, see Herndon, Ash, and Pollin (2014).
On Iceland's alternative path, see Wade and Sigurgeirsd{\'o}ttir (2012).

\section{Exercises}

\begin{exercise}
Construct a doxonomic genealogy of the phrase ``living beyond our means'' in political discourse.
When did it first appear in budget debates?
What institutional channels popularized it?
\end{exercise}

\begin{exercise}
Using the shock doctrine model, analyze the 2008 financial crisis and subsequent austerity turn.
Estimate the narrative entropy trajectory $H(t)$ and identify who supplied the dominant post-crisis frame.
\end{exercise}

\begin{exercise}
Design a survey to measure belief in the household-budget analogy for government spending.
Predict which demographic groups will score highest and explain why.
\end{exercise}
